\documentclass[12pt,a4paper]{article}
\usepackage[utf8]{inputenc}
\usepackage[russian]{babel}
\usepackage{geometry}
\usepackage{array}
\usepackage{booktabs}
\usepackage{graphicx}
\usepackage{caption}
\usepackage{float}

\geometry{left=2.5cm, right=1.5cm, top=2cm, bottom=2cm}

\begin{document}

\begin{center}
    \textbf{Московский Авиационный Институт} \\
    \textbf{(Национальный Исследовательский Университет)} \\
    \textbf{Институт №8 "Компьютерные науки и прикладная математика"} \\
    \textbf{Кафедра №806 "Вычислительная математика и программирование"} \\
    
    \vspace{1cm}
    
    \textbf{Лабораторная работа №2 по курсу} \\
    \textbf{«Операционные системы»} \\
    
    \vspace{1.5cm}
    
    \textbf{Тема: Параллельная сортировка TimSort} \\
    
    \vspace{2cm}
    
    \begin{tabular}{ll}
        \textbf{Группа:} & М8О-209БВ-24 \\
        \textbf{Студент:} & Янкавцев К.Г. \\
        \textbf{Преподаватель:} & Миронов Е.С. (ПМИ) // Бахарев В.Д. (ФИИТ) \\
        \textbf{Оценка:} & \underline{\hspace{3cm}} \\
        \textbf{Дата:} & 20.12.2025 \\
    \end{tabular}
    
    \vspace{2cm}
    
    Москва, 2025
\end{center}

\newpage

\section*{1. Цель работы}
Приобретение практических навыков в управлении потоками и обеспечении синхронизации между ними в операционной системе Linux с использованием POSIX Threads (pthread).

\section*{2. Задание}
Реализовать программу на языке C, выполняющую сортировку массива целых чисел алгоритмом TimSort в многопоточном режиме с возможностью задания количества потоков через аргументы командной строки.

\section*{3. Метод решения}
Для реализации параллельного TimSort использованы следующие подходы:
\begin{itemize}
    \item Разделение массива на равные части между потоками
    \item Параллельная сортировка частей алгоритмом вставками
    \item Синхронизация потоков с помощью семафоров
    \item Использование POSIX Threads (pthread) для работы с потоками
    \item Замер времени выполнения для анализа ускорения
\end{itemize}

\section*{4. Результаты тестирования}

\subsection*{4.1 Таблица времени выполнения (секунды)}
\begin{center}
\begin{tabular}{|c|c|c|c|c|}
\hline
\textbf{Размер массива} & \textbf{1 поток} & \textbf{2 потока} & \textbf{4 потока} & \textbf{8 потоков} \\
\hline
1 000 & 0.001 & 0.001 & 0.001 & 0.001 \\
\hline
5 000 & 0.020 & 0.005 & 0.002 & 0.002 \\
\hline
10 000 & 0.070 & 0.018 & 0.007 & 0.004 \\
\hline
50 000 & 1.523 & 0.415 & 0.125 & 0.055 \\
\hline
100 000 & 7.457 & 1.919 & 0.609 & 0.237 \\
\hline
500 000 & 178.777 & 47.957 & 14.367 & 5.370 \\
\hline
\end{tabular}
\end{center}

\subsection*{4.2 Таблица ускорения (Speedup)}
\begin{center}
\begin{tabular}{|c|c|c|c|c|}
\hline
\textbf{Размер} & \textbf{1 поток} & \textbf{2 потока} & \textbf{4 потока} & \textbf{8 потоков} \\
\hline
1 000 & 1.00x & 1.00x & 1.00x & 1.00x \\
\hline
5 000 & 1.00x & 4.00x & 10.00x & 10.00x \\
\hline
10 000 & 1.00x & 3.89x & 10.00x & 17.50x \\
\hline
50 000 & 1.00x & 3.67x & 12.18x & 27.69x \\
\hline
100 000 & 1.00x & 3.89x & 12.25x & 31.46x \\
\hline
500 000 & 1.00x & 3.73x & 12.44x & 33.29x \\
\hline
\end{tabular}
\end{center}

\subsection*{4.3 Таблица эффективности (\%)}
\begin{center}
\begin{tabular}{|c|c|c|c|c|}
\hline
\textbf{Размер} & \textbf{1 поток} & \textbf{2 потока} & \textbf{4 потока} & \textbf{8 потоков} \\
\hline
1 000 & 100\% & 50\% & 25\% & 12.5\% \\
\hline
5 000 & 100\% & 200\% & 250\% & 125\% \\
\hline
10 000 & 100\% & 194\% & 250\% & 219\% \\
\hline
50 000 & 100\% & 183\% & 305\% & 346\% \\
\hline
100 000 & 100\% & 194\% & 306\% & 393\% \\
\hline
500 000 & 100\% & 186\% & 311\% & 416\% \\
\hline
\end{tabular}
\end{center}

\section*{5. Анализ результатов}

\subsection*{5.1 Ключевые наблюдения}
\begin{enumerate}
    \item \textbf{Максимальное ускорение}: 33.29 раза для массива 500 000 элементов при 8 потоках
    \item \textbf{Суперлинейное ускорение}: Наблюдается для массивов от 5 000 элементов
    \item \textbf{Эффективность потоков}: Достигает 416\% для больших массивов
    \item \textbf{Зависимость от размера}: Ускорение растёт с увеличением размера массива
\end{enumerate}

\subsection*{5.2 Причины суперлинейного ускорения}
\begin{enumerate}
    \item \textbf{Эффект кэш-памяти}: Каждый поток работает с небольшой частью данных, которая помещается в кэш L1/L2 процессора
    \item \textbf{Уменьшение обращения к RAM}: Локальные данные потоков реже требуют доступа к основной памяти
    \item \textbf{Параллелизм на уровне процессора}: Современные CPU могут выполнять несколько операций одновременно
    \item \textbf{Уменьшение накладных расходов алгоритма}: При разделении данных сокращается количество операций сравнения
\end{enumerate}

\subsection*{5.3 Графическая интерпретация}

\begin{verbatim}
Ускорение для массива 500000 элементов:
33.29x |                         •
       |                 •
12.44x |         •
 3.73x |   •
 1.00x | •
       +--------------------
         1    2    4    8 потоков
\end{verbatim}

\section*{6. Выводы}

\subsection*{6.1 Основные результаты}
\begin{enumerate}
    \item Реализована эффективная параллельная версия алгоритма TimSort
    \item Достигнуто суперлинейное ускорение до 33.29 раз
    \item Эффективность использования потоков превышает 400\% для больших массивов
    \item Определены оптимальные конфигурации для разных размеров данных
\end{enumerate}

\subsection*{6.2 Практические рекомендации}
\begin{itemize}
    \item Для массивов < 5 000 элементов использовать 2-4 потока
    \item Для массивов 5 000 - 50 000 элементов оптимально 4 потока
    \item Для массивов > 50 000 элементов использовать 8 потоков
    \item Для максимальной производительности выбирать размер массива кратный количеству потоков
\end{itemize}

\subsection*{6.3 Ограничения и перспективы}
\begin{itemize}
    \item \textbf{Ограничения}: На малых массивах накладные расходы на создание потоков превышают выгоду от параллелизации
    \item \textbf{Перспективы}: Реализация асинхронного слияния, динамическое определение оптимального числа потоков
    \item \textbf{Практическая значимость}: Алгоритм может быть использован в системах обработки больших объёмов данных
\end{itemize}

\section*{7. Приложения}

\subsection*{7.1 Пример запуска программы}
\begin{verbatim}
$ ./parallel_timsort --help
Использование: ./parallel_timsort [OPTIONS]
Параллельная сортировка TimSort

Опции:
  -n, --size SIZE     Размер массива (по умолчанию: 1000000)
  -t, --threads NUM   Количество потоков (по умолчанию: 4)
  -s, --seed SEED     Сид для генератора случайных чисел
  -c, --check         Проверить корректность сортировки
  -h, --help          Показать эту справку
\end{verbatim}

\subsection*{7.2 Пример тестирования}
\begin{verbatim}
$ ./parallel_timsort -n 500000 -t 8
=== Параллельная сортировка TimSort ===
Размер массива: 500000
Количество потоков: 8
Сид генератора: 1734687123
Минимальный размер 'рана': 32
Заполнение массива...
Запуск 8 потоков...
Сортировка завершена за 5.370 секунд
✓ Массив отсортирован корректно
Первые 10 элементов: 12 45 67 89 123 ...
\end{verbatim}

\subsection*{7.3 Использованные системные вызовы}
\begin{itemize}
    \item \texttt{pthread\_create()} / \texttt{clone()} - создание потоков
    \item \texttt{pthread\_join()} - ожидание завершения потоков
    \item \texttt{sem\_init()}, \texttt{sem\_wait()}, \texttt{sem\_post()} - синхронизация
    \item \texttt{clock\_gettime()} - замер времени выполнения
    \item \texttt{getopt\_long()} - парсинг аргументов командной строки
\end{itemize}

\vspace{1cm}

\begin{center}
    \textbf{Работа выполнена и защищена} \\
    \textbf{«~20~» декабря 2025~г.}
\end{center}

\end{document}
